\documentclass[10pt,a4paper]{amsart}
\usepackage{amsmath,amssymb,amsthm,graphicx}
% For commutative diagrams we use Xy-pic
\usepackage[all]{xy}

% We use 2cell for 2-commutative diagrams.
\xyoption{2cell}
% Theorem environments.
%
\theoremstyle{plain}
\newtheorem{theorem}[subsection]{Theorem}
\newtheorem{proposition}[subsection]{Proposition}
\newtheorem{lemma}[subsection]{Lemma}

\theoremstyle{definition}
\newtheorem{definition}[subsection]{Definition}
\newtheorem{example}[subsection]{Example}
\newtheorem{exercise}[subsection]{Exercise}
\newtheorem{situation}[subsection]{Situation}

\theoremstyle{remark}
\newtheorem{remark}[subsection]{Remark}
\newtheorem{remarks}[subsection]{Remarks}

\numberwithin{equation}{subsection}

% Macros
%
\def\lim{\mathop{\mathrm{lim}}\nolimits}
\def\colim{\mathop{\mathrm{colim}}\nolimits}
\def\Spec{\mathop{\mathrm{Spec}}}
\def\Hom{\mathop{\mathrm{Hom}}\nolimits}
\def\Ext{\mathop{\mathrm{Ext}}\nolimits}
\def\SheafHom{\mathop{\mathcal{H}\!\mathit{om}}\nolimits}
\def\SheafExt{\mathop{\mathcal{E}\!\mathit{xt}}\nolimits}
\def\Sch{\mathit{Sch}}
\def\Mor{\mathop{\mathrm{Mor}}\nolimits}
\def\Ob{\mathop{\mathrm{Ob}}\nolimits}
\def\Sh{\mathop{\mathit{Sh}}\nolimits}
\def\NL{\mathop{N\!L}\nolimits}
\def\CH{\mathop{\mathrm{CH}}\nolimits}
\def\proetale{{pro\text{-}\acute{e}tale}}
\def\etale{{\acute{e}tale}}
\def\QCoh{\mathit{QCoh}}
\def\Ker{\mathop{\mathrm{Ker}}}
\def\Im{\mathop{\mathrm{Im}}}
\def\Coker{\mathop{\mathrm{Coker}}}
\def\Coim{\mathop{\mathrm{Coim}}}

% Boxtimes
%
\DeclareMathSymbol{\boxtimes}{\mathbin}{AMSa}{"02}

%
% Macros for moduli stacks/spaces
%
\def\QCohstack{\mathcal{QC}\!\mathit{oh}}
\def\Cohstack{\mathcal{C}\!\mathit{oh}}
\def\Spacesstack{\mathcal{S}\!\mathit{paces}}
\def\Quotfunctor{\mathrm{Quot}}
\def\Hilbfunctor{\mathrm{Hilb}}
\def\Curvesstack{\mathcal{C}\!\mathit{urves}}
\def\Polarizedstack{\mathcal{P}\!\mathit{olarized}}
\def\Complexesstack{\mathcal{C}\!\mathit{omplexes}}
% \Pic is the operator that assigns to X its picard group, usage \Pic(X)
% \Picardstack_{X/B} denotes the Picard stack of X over B
% \Picardfunctor_{X/B} denotes the Picard functor of X over B
\def\Pic{\mathop{\mathrm{Pic}}\nolimits}
\def\Picardstack{\mathcal{P}\!\mathit{ic}}
\def\Picardfunctor{\mathrm{Pic}}
\def\Deformationcategory{\mathcal{D}\!\mathit{ef}}
\newenvironment{reference}{}{}
\newenvironment{slogan}{}{}
\newenvironment{history}{}{}
\usepackage{cleveref}
\crefname{originallemma}{original lemma}{original lemmas}
\Crefname{originallemma}{Original lemma}{Original lemmas}
\makeatletter
\newlabel{more-morphisms-lemma-flat-lci}{{069K}{source}{}{sourcepointer-1}{}}
\newlabel{more-morphisms-lemma-flat-lci@cref}{{[originallemma][0][]069K}{[1][0][]source}{}{}{}}
\makeatother
\crefname{originallemma}{original lemma}{original lemmas}
\Crefname{originallemma}{Original lemma}{Original lemmas}
\makeatletter
\newlabel{morphisms-lemma-syntomic-flat-fibres}{{01UF}{source}{}{sourcepointer-2}{}}
\newlabel{morphisms-lemma-syntomic-flat-fibres@cref}{{[originallemma][0][]01UF}{[1][0][]source}{}{}{}}
\makeatother
\crefname{originaldefinition}{original definition}{original definitions}
\Crefname{originaldefinition}{Original definition}{Original definitions}
\makeatletter
\newlabel{algebra-definition-relative-global-complete-intersection}{{00SP}{source}{}{sourcepointer-3}{}}
\newlabel{algebra-definition-relative-global-complete-intersection@cref}{{[originaldefinition][0][]00SP}{[1][0][]source}{}{}{}}
\makeatother
\crefname{originallemma}{original lemma}{original lemmas}
\Crefname{originallemma}{Original lemma}{Original lemmas}
\makeatletter
\newlabel{algebra-lemma-relative-global-complete-intersection}{{00SW}{source}{}{sourcepointer-4}{}}
\newlabel{algebra-lemma-relative-global-complete-intersection@cref}{{[originallemma][0][]00SW}{[1][0][]source}{}{}{}}
\makeatother
\crefname{originallemma}{original lemma}{original lemmas}
\Crefname{originallemma}{Original lemma}{Original lemmas}
\makeatletter
\newlabel{algebra-lemma-isolated-point-fibre}{{00PK}{source}{}{sourcepointer-5}{}}
\newlabel{algebra-lemma-isolated-point-fibre@cref}{{[originallemma][0][]00PK}{[1][0][]source}{}{}{}}
\makeatother
\crefname{originallemma}{original lemma}{original lemmas}
\Crefname{originallemma}{Original lemma}{Original lemmas}
\makeatletter
\newlabel{algebra-lemma-syntomic}{{00SY}{source}{}{sourcepointer-6}{}}
\newlabel{algebra-lemma-syntomic@cref}{{[originallemma][0][]00SY}{[1][0][]source}{}{}{}}
\makeatother
\crefname{originallemma}{original lemma}{original lemmas}
\Crefname{originallemma}{Original lemma}{Original lemmas}
\makeatletter
\newlabel{morphisms-lemma-syntomic-locally-standard-syntomic}{{01UE}{source}{}{sourcepointer-7}{}}
\newlabel{morphisms-lemma-syntomic-locally-standard-syntomic@cref}{{[originallemma][0][]01UE}{[1][0][]source}{}{}{}}
\makeatother
\crefname{originallemma}{original lemma}{original lemmas}
\Crefname{originallemma}{Original lemma}{Original lemmas}
\makeatletter
\newlabel{more-morphisms-lemma-flat-fp-NL-lci}{{0FK3}{source}{}{sourcepointer-8}{}}
\newlabel{more-morphisms-lemma-flat-fp-NL-lci@cref}{{[originallemma][0][]0FK3}{[1][0][]source}{}{}{}}
\makeatother
\crefname{originallemma}{original lemma}{original lemmas}
\Crefname{originallemma}{Original lemma}{Original lemmas}
\makeatletter
\newlabel{algebra-lemma-relative-global-complete-intersection-conormal}{{00SV}{source}{}{sourcepointer-9}{}}
\newlabel{algebra-lemma-relative-global-complete-intersection-conormal@cref}{{[originallemma][0][]00SV}{[1][0][]source}{}{}{}}
\makeatother
\crefname{originallemma}{original lemma}{original lemmas}
\Crefname{originallemma}{Original lemma}{Original lemmas}
\makeatletter
\newlabel{lemma-dualizing-base-flat-flat}{{0BVE}{source}{}{sourcepointer-10}{}}
\newlabel{lemma-dualizing-base-flat-flat@cref}{{[originallemma][0][]0BVE}{[1][0][]source}{}{}{}}
\makeatother
\crefname{originallemma}{original lemma}{original lemmas}
\Crefname{originallemma}{Original lemma}{Original lemmas}
\makeatletter
\newlabel{more-morphisms-lemma-flat-and-free-at-point-fibre}{{080Q}{source}{}{sourcepointer-11}{}}
\newlabel{more-morphisms-lemma-flat-and-free-at-point-fibre@cref}{{[originallemma][0][]080Q}{[1][0][]source}{}{}{}}
\makeatother
\crefname{originallemma}{original lemma}{original lemmas}
\Crefname{originallemma}{Original lemma}{Original lemmas}
\makeatletter
\newlabel{lemma-dualizing-base-change-of-flat}{{0BVF}{source}{}{sourcepointer-12}{}}
\newlabel{lemma-dualizing-base-change-of-flat@cref}{{[originallemma][0][]0BVF}{[1][0][]source}{}{}{}}
\makeatother
\crefname{originalsection}{original section}{original sections}
\Crefname{originalsection}{Original section}{Original sections}
\makeatletter
\newlabel{remark-relative-dualizing-for-quasi-finite}{{0BVG}{source}{}{sourcepointer-13}{}}
\newlabel{remark-relative-dualizing-for-quasi-finite@cref}{{[originalsection][0][]0BVG}{[1][0][]source}{}{}{}}
\makeatother
\crefname{originallemma}{original lemma}{original lemmas}
\Crefname{originallemma}{Original lemma}{Original lemmas}
\makeatletter
\newlabel{algebra-lemma-Noether-normalization}{{00OY}{source}{}{sourcepointer-14}{}}
\newlabel{algebra-lemma-Noether-normalization@cref}{{[originallemma][0][]00OY}{[1][0][]source}{}{}{}}
\makeatother
\crefname{originallemma}{original lemma}{original lemmas}
\Crefname{originallemma}{Original lemma}{Original lemmas}
\makeatletter
\newlabel{dualizing-lemma-dualizing-finite}{{0AX0}{source}{}{sourcepointer-15}{}}
\newlabel{dualizing-lemma-dualizing-finite@cref}{{[originallemma][0][]0AX0}{[1][0][]source}{}{}{}}
\makeatother
\crefname{originallemma}{original lemma}{original lemmas}
\Crefname{originallemma}{Original lemma}{Original lemmas}
\makeatletter
\newlabel{dualizing-lemma-RHom-ext}{{0A71}{source}{}{sourcepointer-16}{}}
\newlabel{dualizing-lemma-RHom-ext@cref}{{[originallemma][0][]0A71}{[1][0][]source}{}{}{}}
\makeatother
\crefname{originallemma}{original lemma}{original lemmas}
\Crefname{originallemma}{Original lemma}{Original lemmas}
\makeatletter
\newlabel{dualizing-lemma-gorenstein}{{0DW9}{source}{}{sourcepointer-17}{}}
\newlabel{dualizing-lemma-gorenstein@cref}{{[originallemma][0][]0DW9}{[1][0][]source}{}{}{}}
\makeatother
\crefname{originallemma}{original lemma}{original lemmas}
\Crefname{originallemma}{Original lemma}{Original lemmas}
\makeatletter
\newlabel{dualizing-lemma-gorenstein-lci}{{0DWA}{source}{}{sourcepointer-18}{}}
\newlabel{dualizing-lemma-gorenstein-lci@cref}{{[originallemma][0][]0DWA}{[1][0][]source}{}{}{}}
\makeatother
\crefname{originallemma}{original lemma}{original lemmas}
\Crefname{originallemma}{Original lemma}{Original lemmas}
\makeatletter
\newlabel{algebra-lemma-quasi-finite-open}{{00QA}{source}{}{sourcepointer-19}{}}
\newlabel{algebra-lemma-quasi-finite-open@cref}{{[originallemma][0][]00QA}{[1][0][]source}{}{}{}}
\makeatother
\crefname{originallemma}{original lemma}{original lemmas}
\Crefname{originallemma}{Original lemma}{Original lemmas}
\makeatletter
\newlabel{morphisms-lemma-etale-at-point}{{02GU}{source}{}{sourcepointer-20}{}}
\newlabel{morphisms-lemma-etale-at-point@cref}{{[originallemma][0][]02GU}{[1][0][]source}{}{}{}}
\makeatother
\crefname{originallemma}{original lemma}{original lemmas}
\Crefname{originallemma}{Original lemma}{Original lemmas}
\makeatletter
\newlabel{more-algebra-lemma-functorial}{{0624}{source}{}{sourcepointer-21}{}}
\newlabel{more-algebra-lemma-functorial@cref}{{[originallemma][0][]0624}{[1][0][]source}{}{}{}}
\makeatother
\crefname{originallemma}{original lemma}{original lemmas}
\Crefname{originallemma}{Original lemma}{Original lemmas}
\makeatletter
\newlabel{more-algebra-lemma-relative-global-complete-intersection-koszul}{{07D2}{source}{}{sourcepointer-22}{}}
\newlabel{more-algebra-lemma-relative-global-complete-intersection-koszul@cref}{{[originallemma][0][]07D2}{[1][0][]source}{}{}{}}
\makeatother
\crefname{originallemma}{original lemma}{original lemmas}
\Crefname{originallemma}{Original lemma}{Original lemmas}
\makeatletter
\newlabel{more-algebra-lemma-truncate-koszul-regular}{{068M}{source}{}{sourcepointer-23}{}}
\newlabel{more-algebra-lemma-truncate-koszul-regular@cref}{{[originallemma][0][]068M}{[1][0][]source}{}{}{}}
\makeatother
\crefname{originallemma}{original lemma}{original lemmas}
\Crefname{originallemma}{Original lemma}{Original lemmas}
\makeatletter
\newlabel{more-algebra-lemma-koszul-regular-flat-base-change}{{062H}{source}{}{sourcepointer-24}{}}
\newlabel{more-algebra-lemma-koszul-regular-flat-base-change@cref}{{[originallemma][0][]062H}{[1][0][]source}{}{}{}}
\makeatother
\crefname{originallemma}{original lemma}{original lemmas}
\Crefname{originallemma}{Original lemma}{Original lemmas}
\makeatletter
\newlabel{lemma-flat-gorenstein-agree-noether}{{0BW6}{source}{}{sourcepointer-25}{}}
\newlabel{lemma-flat-gorenstein-agree-noether@cref}{{[originallemma][0][]0BW6}{[1][0][]source}{}{}{}}
\makeatother
\crefname{originallemma}{original lemma}{original lemmas}
\Crefname{originallemma}{Original lemma}{Original lemmas}
\makeatletter
\newlabel{lemma-kahler-different-complete-intersection}{{0BVZ}{source}{}{sourcepointer-26}{}}
\newlabel{lemma-kahler-different-complete-intersection@cref}{{[originallemma][0][]0BVZ}{[1][0][]source}{}{}{}}
\makeatother
\crefname{originallemma}{original lemma}{original lemmas}
\Crefname{originallemma}{Original lemma}{Original lemmas}
\makeatletter
\newlabel{lemma-noether-different-flat-quasi-finite}{{0BVT}{source}{}{sourcepointer-27}{}}
\newlabel{lemma-noether-different-flat-quasi-finite@cref}{{[originallemma][0][]0BVT}{[1][0][]source}{}{}{}}
\makeatother
\crefname{originallemma}{original lemma}{original lemmas}
\Crefname{originallemma}{Original lemma}{Original lemmas}
\makeatletter
\newlabel{lemma-different-ramification}{{0BW9}{source}{}{sourcepointer-28}{}}
\newlabel{lemma-different-ramification@cref}{{[originallemma][0][]0BW9}{[1][0][]source}{}{}{}}
\makeatother
\crefname{originallemma}{original lemma}{original lemmas}
\Crefname{originallemma}{Original lemma}{Original lemmas}
\makeatletter
\newlabel{properties-lemma-standard-open-two-affines}{{0H9B}{source}{}{sourcepointer-29}{}}
\newlabel{properties-lemma-standard-open-two-affines@cref}{{[originallemma][0][]0H9B}{[1][0][]source}{}{}{}}
\makeatother
\crefname{originallemma}{original lemma}{original lemmas}
\Crefname{originallemma}{Original lemma}{Original lemmas}
\makeatletter
\newlabel{more-algebra-lemma-module-over-fibre-product}{{07RU}{source}{}{sourcepointer-30}{}}
\newlabel{more-algebra-lemma-module-over-fibre-product@cref}{{[originallemma][0][]07RU}{[1][0][]source}{}{}{}}
\makeatother
\crefname{originallemma}{original lemma}{original lemmas}
\Crefname{originallemma}{Original lemma}{Original lemmas}
\makeatletter
\newlabel{more-algebra-lemma-properties-algebras-over-fibre-product}{{08KQ}{source}{}{sourcepointer-31}{}}
\newlabel{more-algebra-lemma-properties-algebras-over-fibre-product@cref}{{[originallemma][0][]08KQ}{[1][0][]source}{}{}{}}
\makeatother
\crefname{originallemma}{original lemma}{original lemmas}
\Crefname{originallemma}{Original lemma}{Original lemmas}
\makeatletter
\newlabel{limits-lemma-descend-finite-presentation}{{01ZM}{source}{}{sourcepointer-32}{}}
\newlabel{limits-lemma-descend-finite-presentation@cref}{{[originallemma][0][]01ZM}{[1][0][]source}{}{}{}}
\makeatother
\crefname{originallemma}{original lemma}{original lemmas}
\Crefname{originallemma}{Original lemma}{Original lemmas}
\makeatletter
\newlabel{limits-lemma-descend-syntomic}{{0C3L}{source}{}{sourcepointer-33}{}}
\newlabel{limits-lemma-descend-syntomic@cref}{{[originallemma][0][]0C3L}{[1][0][]source}{}{}{}}
\makeatother
\crefname{originallemma}{original lemma}{original lemmas}
\Crefname{originallemma}{Original lemma}{Original lemmas}
\makeatletter
\newlabel{limits-lemma-descend-quasi-finite}{{094M}{source}{}{sourcepointer-34}{}}
\newlabel{limits-lemma-descend-quasi-finite@cref}{{[originallemma][0][]094M}{[1][0][]source}{}{}{}}
\makeatother
\crefname{originallemma}{original lemma}{original lemmas}
\Crefname{originallemma}{Original lemma}{Original lemmas}
\makeatletter
\newlabel{perfect-lemma-determinant-two-term-complexes}{{0FJX}{source}{}{sourcepointer-35}{}}
\newlabel{perfect-lemma-determinant-two-term-complexes@cref}{{[originallemma][0][]0FJX}{[1][0][]source}{}{}{}}
\makeatother
\crefname{originallemma}{original lemma}{original lemmas}
\Crefname{originallemma}{Original lemma}{Original lemmas}
\makeatletter
\newlabel{more-morphisms-lemma-base-change-NL-flat}{{0FK0}{source}{}{sourcepointer-36}{}}
\newlabel{more-morphisms-lemma-base-change-NL-flat@cref}{{[originallemma][0][]0FK0}{[1][0][]source}{}{}{}}
\makeatother
\crefname{originalsection}{original section}{original sections}
\Crefname{originalsection}{Original section}{Original sections}
\makeatletter
\newlabel{perfect-remark-functorial-det}{{0FJY}{source}{}{sourcepointer-37}{}}
\newlabel{perfect-remark-functorial-det@cref}{{[originalsection][0][]0FJY}{[1][0][]source}{}{}{}}
\makeatother
\crefname{originalsection}{original section}{original sections}
\Crefname{originalsection}{Original section}{Original sections}
\makeatletter
\newlabel{remark-relative-dualizing-for-flat-quasi-finite}{{0BVJ}{source}{}{sourcepointer-38}{}}
\newlabel{remark-relative-dualizing-for-flat-quasi-finite@cref}{{[originalsection][0][]0BVJ}{[1][0][]source}{}{}{}}
\makeatother
\crefname{originalequation}{original equation}{original equations}
\Crefname{originalequation}{Original equation}{Original equations}
\makeatletter
\newlabel{equation-bc-dualizing}{{0BVB}{source}{}{sourcepointer-39}{}}
\newlabel{equation-bc-dualizing@cref}{{[originalequation][0][]0BVB}{[1][0][]source}{}{}{}}
\makeatother
\crefname{originallemma}{original lemma}{original lemmas}
\Crefname{originallemma}{Original lemma}{Original lemmas}
\makeatletter
\newlabel{lemma-trace-base-change}{{0BT9}{source}{}{sourcepointer-40}{}}
\newlabel{lemma-trace-base-change@cref}{{[originallemma][0][]0BT9}{[1][0][]source}{}{}{}}
\makeatother
\title{Canonical Tate determinant-to-dualizing isomorphism for quasi-finite syntomic maps}
\author{The Stacks Project Authors}
\date{stacks-project:discriminant.tex:local-snapshot-20261011; GNU FDL 1.2 or later}
\begin{document}
\maketitle
\noindent\textit{Editorial note.} In the fibre-product display the second tensor product is over A\_2, as fixed by both its factors and the following base-change identity. The printed A\_1 subscript is retained and disclosed.\par\medskip
\noindent\textit{Editorial note.} In the Noetherian approximation paragraph, the finite polynomial subalgebra notation is missing a closing bracket; no mathematical variable or hypothesis is changed.\par\medskip
\noindent\textit{Editorial note.} A Garel differential-form construction in the adjacent chapter uses a related descent pattern but a different map with additional exterior-form obligations; it is not a second proof of this determinant/dualizing isomorphism.\par\medskip
\noindent\textit{Allowed standard prerequisites (editorial).}\begin{itemize}
\item Explicit prerequisite boundary. Determinant of the naive cotangent complex, dualizing modules for quasi-finite maps, trace base change, flat finite-presentation patching over fibre products, and finite-presentation limit descent. The local different formula and universal syntomic family are retained author support.
\end{itemize}\medskip
\ifdefined\hypertarget\hypertarget{sourcepointer-1}{}\fi
\noindent\textit{Source reference.} Original lemma 069K: focused-seeds/AGP/sources/stacks-snapshot-20261011/tags.txt, source lines 8044--8044. Stable Stacks tag for an expressly allowed foundational or notational input; the counted construction remains the two concrete core obligations. The retained local core is not replaced by this pointer. See standard\_prerequisites.\par\medskip
\ifdefined\hypertarget\hypertarget{sourcepointer-2}{}\fi
\noindent\textit{Source reference.} Original lemma 01UF: focused-seeds/AGP/sources/stacks-snapshot-20261011/tags.txt, source lines 2408--2408. Stable Stacks tag for an expressly allowed foundational or notational input; the counted construction remains the two concrete core obligations. The retained local core is not replaced by this pointer. See standard\_prerequisites.\par\medskip
\ifdefined\hypertarget\hypertarget{sourcepointer-3}{}\fi
\noindent\textit{Source reference.} Original definition 00SP: focused-seeds/AGP/sources/stacks-snapshot-20261011/tags.txt, source lines 1047--1047. Stable Stacks tag for an expressly allowed foundational or notational input; the counted construction remains the two concrete core obligations. The retained local core is not replaced by this pointer. See standard\_prerequisites.\par\medskip
\ifdefined\hypertarget\hypertarget{sourcepointer-4}{}\fi
\noindent\textit{Source reference.} Original lemma 00SW: focused-seeds/AGP/sources/stacks-snapshot-20261011/tags.txt, source lines 1054--1054. Stable Stacks tag for an expressly allowed foundational or notational input; the counted construction remains the two concrete core obligations. The retained local core is not replaced by this pointer. See standard\_prerequisites.\par\medskip
\ifdefined\hypertarget\hypertarget{sourcepointer-5}{}\fi
\noindent\textit{Source reference.} Original lemma 00PK: focused-seeds/AGP/sources/stacks-snapshot-20261011/tags.txt, source lines 934--934. Stable Stacks tag for an expressly allowed foundational or notational input; the counted construction remains the two concrete core obligations. The retained local core is not replaced by this pointer. See standard\_prerequisites.\par\medskip
\ifdefined\hypertarget\hypertarget{sourcepointer-6}{}\fi
\noindent\textit{Source reference.} Original lemma 00SY: focused-seeds/AGP/sources/stacks-snapshot-20261011/tags.txt, source lines 1056--1056. Stable Stacks tag for an expressly allowed foundational or notational input; the counted construction remains the two concrete core obligations. The retained local core is not replaced by this pointer. See standard\_prerequisites.\par\medskip
\ifdefined\hypertarget\hypertarget{sourcepointer-7}{}\fi
\noindent\textit{Source reference.} Original lemma 01UE: focused-seeds/AGP/sources/stacks-snapshot-20261011/tags.txt, source lines 2407--2407. Stable Stacks tag for an expressly allowed foundational or notational input; the counted construction remains the two concrete core obligations. The retained local core is not replaced by this pointer. See standard\_prerequisites.\par\medskip
\ifdefined\hypertarget\hypertarget{sourcepointer-8}{}\fi
\noindent\textit{Source reference.} Original lemma 0FK3: focused-seeds/AGP/sources/stacks-snapshot-20261011/tags.txt, source lines 19448--19448. Stable Stacks tag for an expressly allowed foundational or notational input; the counted construction remains the two concrete core obligations. The retained local core is not replaced by this pointer. See standard\_prerequisites.\par\medskip
\ifdefined\hypertarget\hypertarget{sourcepointer-9}{}\fi
\noindent\textit{Source reference.} Original lemma 00SV: focused-seeds/AGP/sources/stacks-snapshot-20261011/tags.txt, source lines 1053--1053. Stable Stacks tag for an expressly allowed foundational or notational input; the counted construction remains the two concrete core obligations. The retained local core is not replaced by this pointer. See standard\_prerequisites.\par\medskip
\ifdefined\hypertarget\hypertarget{sourcepointer-10}{}\fi
\noindent\textit{Source reference.} Original lemma 0BVE: focused-seeds/AGP/sources/stacks-snapshot-20261011/tags.txt, source lines 14907--14907. Stable Stacks tag for an expressly allowed foundational or notational input; the counted construction remains the two concrete core obligations. The retained local core is not replaced by this pointer. See standard\_prerequisites.\par\medskip
\ifdefined\hypertarget\hypertarget{sourcepointer-11}{}\fi
\noindent\textit{Source reference.} Original lemma 080Q: focused-seeds/AGP/sources/stacks-snapshot-20261011/tags.txt, source lines 10190--10190. Stable Stacks tag for an expressly allowed foundational or notational input; the counted construction remains the two concrete core obligations. The retained local core is not replaced by this pointer. See standard\_prerequisites.\par\medskip
\ifdefined\hypertarget\hypertarget{sourcepointer-12}{}\fi
\noindent\textit{Source reference.} Original lemma 0BVF: focused-seeds/AGP/sources/stacks-snapshot-20261011/tags.txt, source lines 14908--14908. Stable Stacks tag for an expressly allowed foundational or notational input; the counted construction remains the two concrete core obligations. The retained local core is not replaced by this pointer. See standard\_prerequisites.\par\medskip
\ifdefined\hypertarget\hypertarget{sourcepointer-13}{}\fi
\noindent\textit{Source reference.} Original section 0BVG: focused-seeds/AGP/sources/stacks-snapshot-20261011/tags.txt, source lines 14909--14909. Stable Stacks tag for an expressly allowed foundational or notational input; the counted construction remains the two concrete core obligations. The retained local core is not replaced by this pointer. See standard\_prerequisites.\par\medskip
\ifdefined\hypertarget\hypertarget{sourcepointer-14}{}\fi
\noindent\textit{Source reference.} Original lemma 00OY: focused-seeds/AGP/sources/stacks-snapshot-20261011/tags.txt, source lines 912--912. Stable Stacks tag for an expressly allowed foundational or notational input; the counted construction remains the two concrete core obligations. The retained local core is not replaced by this pointer. See standard\_prerequisites.\par\medskip
\ifdefined\hypertarget\hypertarget{sourcepointer-15}{}\fi
\noindent\textit{Source reference.} Original lemma 0AX0: focused-seeds/AGP/sources/stacks-snapshot-20261011/tags.txt, source lines 13738--13738. Stable Stacks tag for an expressly allowed foundational or notational input; the counted construction remains the two concrete core obligations. The retained local core is not replaced by this pointer. See standard\_prerequisites.\par\medskip
\ifdefined\hypertarget\hypertarget{sourcepointer-16}{}\fi
\noindent\textit{Source reference.} Original lemma 0A71: focused-seeds/AGP/sources/stacks-snapshot-20261011/tags.txt, source lines 12864--12864. Stable Stacks tag for an expressly allowed foundational or notational input; the counted construction remains the two concrete core obligations. The retained local core is not replaced by this pointer. See standard\_prerequisites.\par\medskip
\ifdefined\hypertarget\hypertarget{sourcepointer-17}{}\fi
\noindent\textit{Source reference.} Original lemma 0DW9: focused-seeds/AGP/sources/stacks-snapshot-20261011/tags.txt, source lines 17389--17389. Stable Stacks tag for an expressly allowed foundational or notational input; the counted construction remains the two concrete core obligations. The retained local core is not replaced by this pointer. See standard\_prerequisites.\par\medskip
\ifdefined\hypertarget\hypertarget{sourcepointer-18}{}\fi
\noindent\textit{Source reference.} Original lemma 0DWA: focused-seeds/AGP/sources/stacks-snapshot-20261011/tags.txt, source lines 17390--17390. Stable Stacks tag for an expressly allowed foundational or notational input; the counted construction remains the two concrete core obligations. The retained local core is not replaced by this pointer. See standard\_prerequisites.\par\medskip
\ifdefined\hypertarget\hypertarget{sourcepointer-19}{}\fi
\noindent\textit{Source reference.} Original lemma 00QA: focused-seeds/AGP/sources/stacks-snapshot-20261011/tags.txt, source lines 960--960. Stable Stacks tag for an expressly allowed foundational or notational input; the counted construction remains the two concrete core obligations. The retained local core is not replaced by this pointer. See standard\_prerequisites.\par\medskip
\ifdefined\hypertarget\hypertarget{sourcepointer-20}{}\fi
\noindent\textit{Source reference.} Original lemma 02GU: focused-seeds/AGP/sources/stacks-snapshot-20261011/tags.txt, source lines 3241--3241. Stable Stacks tag for an expressly allowed foundational or notational input; the counted construction remains the two concrete core obligations. The retained local core is not replaced by this pointer. See standard\_prerequisites.\par\medskip
\ifdefined\hypertarget\hypertarget{sourcepointer-21}{}\fi
\noindent\textit{Source reference.} Original lemma 0624: focused-seeds/AGP/sources/stacks-snapshot-20261011/tags.txt, source lines 7783--7783. Stable Stacks tag for an expressly allowed foundational or notational input; the counted construction remains the two concrete core obligations. The retained local core is not replaced by this pointer. See standard\_prerequisites.\par\medskip
\ifdefined\hypertarget\hypertarget{sourcepointer-22}{}\fi
\noindent\textit{Source reference.} Original lemma 07D2: focused-seeds/AGP/sources/stacks-snapshot-20261011/tags.txt, source lines 9397--9397. Stable Stacks tag for an expressly allowed foundational or notational input; the counted construction remains the two concrete core obligations. The retained local core is not replaced by this pointer. See standard\_prerequisites.\par\medskip
\ifdefined\hypertarget\hypertarget{sourcepointer-23}{}\fi
\noindent\textit{Source reference.} Original lemma 068M: focused-seeds/AGP/sources/stacks-snapshot-20261011/tags.txt, source lines 8011--8011. Stable Stacks tag for an expressly allowed foundational or notational input; the counted construction remains the two concrete core obligations. The retained local core is not replaced by this pointer. See standard\_prerequisites.\par\medskip
\ifdefined\hypertarget\hypertarget{sourcepointer-24}{}\fi
\noindent\textit{Source reference.} Original lemma 062H: focused-seeds/AGP/sources/stacks-snapshot-20261011/tags.txt, source lines 7796--7796. Stable Stacks tag for an expressly allowed foundational or notational input; the counted construction remains the two concrete core obligations. The retained local core is not replaced by this pointer. See standard\_prerequisites.\par\medskip
\ifdefined\hypertarget\hypertarget{sourcepointer-25}{}\fi
\noindent\textit{Source reference.} Original lemma 0BW6: focused-seeds/AGP/sources/stacks-snapshot-20261011/tags.txt, source lines 14934--14934. Stable Stacks tag for an expressly allowed foundational or notational input; the counted construction remains the two concrete core obligations. The retained local core is not replaced by this pointer. See standard\_prerequisites.\par\medskip
\ifdefined\hypertarget\hypertarget{sourcepointer-26}{}\fi
\noindent\textit{Source reference.} Original lemma 0BVZ: focused-seeds/AGP/sources/stacks-snapshot-20261011/tags.txt, source lines 14927--14927. Stable Stacks tag for an expressly allowed foundational or notational input; the counted construction remains the two concrete core obligations. The retained local core is not replaced by this pointer. See standard\_prerequisites.\par\medskip
\ifdefined\hypertarget\hypertarget{sourcepointer-27}{}\fi
\noindent\textit{Source reference.} Original lemma 0BVT: focused-seeds/AGP/sources/stacks-snapshot-20261011/tags.txt, source lines 14921--14921. Stable Stacks tag for an expressly allowed foundational or notational input; the counted construction remains the two concrete core obligations. The retained local core is not replaced by this pointer. See standard\_prerequisites.\par\medskip
\ifdefined\hypertarget\hypertarget{sourcepointer-28}{}\fi
\noindent\textit{Source reference.} Original lemma 0BW9: focused-seeds/AGP/sources/stacks-snapshot-20261011/tags.txt, source lines 14937--14937. Stable Stacks tag for an expressly allowed foundational or notational input; the counted construction remains the two concrete core obligations. The retained local core is not replaced by this pointer. See standard\_prerequisites.\par\medskip
\ifdefined\hypertarget\hypertarget{sourcepointer-29}{}\fi
\noindent\textit{Source reference.} Original lemma 0H9B: focused-seeds/AGP/sources/stacks-snapshot-20261011/tags.txt, source lines 21521--21521. Stable Stacks tag for an expressly allowed foundational or notational input; the counted construction remains the two concrete core obligations. The retained local core is not replaced by this pointer. See standard\_prerequisites.\par\medskip
\ifdefined\hypertarget\hypertarget{sourcepointer-30}{}\fi
\noindent\textit{Source reference.} Original lemma 07RU: focused-seeds/AGP/sources/stacks-snapshot-20261011/tags.txt, source lines 9879--9879. Stable Stacks tag for an expressly allowed foundational or notational input; the counted construction remains the two concrete core obligations. The retained local core is not replaced by this pointer. See standard\_prerequisites.\par\medskip
\ifdefined\hypertarget\hypertarget{sourcepointer-31}{}\fi
\noindent\textit{Source reference.} Original lemma 08KQ: focused-seeds/AGP/sources/stacks-snapshot-20261011/tags.txt, source lines 10893--10893. Stable Stacks tag for an expressly allowed foundational or notational input; the counted construction remains the two concrete core obligations. The retained local core is not replaced by this pointer. See standard\_prerequisites.\par\medskip
\ifdefined\hypertarget\hypertarget{sourcepointer-32}{}\fi
\noindent\textit{Source reference.} Original lemma 01ZM: focused-seeds/AGP/sources/stacks-snapshot-20261011/tags.txt, source lines 2599--2599. Stable Stacks tag for an expressly allowed foundational or notational input; the counted construction remains the two concrete core obligations. The retained local core is not replaced by this pointer. See standard\_prerequisites.\par\medskip
\ifdefined\hypertarget\hypertarget{sourcepointer-33}{}\fi
\noindent\textit{Source reference.} Original lemma 0C3L: focused-seeds/AGP/sources/stacks-snapshot-20261011/tags.txt, source lines 15196--15196. Stable Stacks tag for an expressly allowed foundational or notational input; the counted construction remains the two concrete core obligations. The retained local core is not replaced by this pointer. See standard\_prerequisites.\par\medskip
\ifdefined\hypertarget\hypertarget{sourcepointer-34}{}\fi
\noindent\textit{Source reference.} Original lemma 094M: focused-seeds/AGP/sources/stacks-snapshot-20261011/tags.txt, source lines 11555--11555. Stable Stacks tag for an expressly allowed foundational or notational input; the counted construction remains the two concrete core obligations. The retained local core is not replaced by this pointer. See standard\_prerequisites.\par\medskip
\ifdefined\hypertarget\hypertarget{sourcepointer-35}{}\fi
\noindent\textit{Source reference.} Original lemma 0FJX: focused-seeds/AGP/sources/stacks-snapshot-20261011/tags.txt, source lines 19442--19442. Stable Stacks tag for an expressly allowed foundational or notational input; the counted construction remains the two concrete core obligations. The retained local core is not replaced by this pointer. See standard\_prerequisites.\par\medskip
\ifdefined\hypertarget\hypertarget{sourcepointer-36}{}\fi
\noindent\textit{Source reference.} Original lemma 0FK0: focused-seeds/AGP/sources/stacks-snapshot-20261011/tags.txt, source lines 19445--19445. Stable Stacks tag for an expressly allowed foundational or notational input; the counted construction remains the two concrete core obligations. The retained local core is not replaced by this pointer. See standard\_prerequisites.\par\medskip
\ifdefined\hypertarget\hypertarget{sourcepointer-37}{}\fi
\noindent\textit{Source reference.} Original section 0FJY: focused-seeds/AGP/sources/stacks-snapshot-20261011/tags.txt, source lines 19443--19443. Stable Stacks tag for an expressly allowed foundational or notational input; the counted construction remains the two concrete core obligations. The retained local core is not replaced by this pointer. See standard\_prerequisites.\par\medskip
\ifdefined\hypertarget\hypertarget{sourcepointer-38}{}\fi
\noindent\textit{Source reference.} Original section 0BVJ: focused-seeds/AGP/sources/stacks-snapshot-20261011/tags.txt, source lines 14912--14912. Stable Stacks tag for an expressly allowed foundational or notational input; the counted construction remains the two concrete core obligations. The retained local core is not replaced by this pointer. See standard\_prerequisites.\par\medskip
\ifdefined\hypertarget\hypertarget{sourcepointer-39}{}\fi
\noindent\textit{Source reference.} Original equation 0BVB: focused-seeds/AGP/sources/stacks-snapshot-20261011/tags.txt, source lines 14904--14904. Stable Stacks tag for an expressly allowed foundational or notational input; the counted construction remains the two concrete core obligations. The retained local core is not replaced by this pointer. See standard\_prerequisites.\par\medskip
\ifdefined\hypertarget\hypertarget{sourcepointer-40}{}\fi
\noindent\textit{Source reference.} Original lemma 0BT9: focused-seeds/AGP/sources/stacks-snapshot-20261011/tags.txt, source lines 14832--14832. Stable Stacks tag for an expressly allowed foundational or notational input; the counted construction remains the two concrete core obligations. The retained local core is not replaced by this pointer. See standard\_prerequisites.\par\medskip
\begin{proposition}
\label{proposition-tate-map}
There exists a unique rule that to every locally quasi-finite syntomic
morphism of locally Noetherian schemes $Y \to X$ assigns an isomorphism
$$
c_{Y/X} : \det(\NL_{Y/X}) \longrightarrow \omega_{Y/X}
$$
satisfying the following two properties
\begin{enumerate}
\item the section $\delta(\NL_{Y/X})$ is mapped to $\tau_{Y/X}$, and
\item the rule is compatible with restriction to opens and with
base change.
\end{enumerate}
\end{proposition}

\section{The Tate map}
\label{section-tate-map}

\noindent
In this section we produce an isomorphism between
the determinant of the relative cotangent complex and
the relative dualizing module for a locally quasi-finite
syntomic morphism of locally Noetherian schemes. Following
\cite[1.4.4]{Garel} we dub the isomorphism the Tate map.
Our approach is to avoid doing local calculations as
much as is possible.

\medskip\noindent
Let $Y \to X$ be a locally quasi-finite syntomic morphism of schemes.
We will use all the equivalent conditions for this notion given in
Lemma \ref{lemma-syntomic-quasi-finite} without further mention in
this section. In particular, we see that $\NL_{Y/X}$ is a perfect
object of $D(\mathcal{O}_Y)$ with tor-amplitude in $[-1, 0]$. Thus
we have a canonical invertible module
$\det(\NL_{Y/X})$ on $Y$ and a global section
$$
\delta(\NL_{Y/X}) \in \Gamma(Y, \det(\NL_{Y/X}))
$$
See Derived Categories of Schemes, Lemma
\ref{perfect-lemma-determinant-two-term-complexes}.
Suppose given a commutative diagram of schemes
$$
\xymatrix{
Y' \ar[r]_b \ar[d] & Y \ar[d] \\
X' \ar[r] & X
}
$$
whose vertical arrows are locally quasi-finite syntomic and which
induces an isomorphism of $Y'$ with an open of $X' \times_X Y$.
Then the canonical map
$$
Lb^*\NL_{Y/X} \longrightarrow \NL_{Y'/X'}
$$
is a quasi-isomorphism by
More on Morphisms, Lemma \ref{more-morphisms-lemma-base-change-NL-flat}.
Thus we get a canonical isomorphism
$b^*\det(\NL_{Y/X}) \to \det(\NL_{Y'/X'})$ which sends the
canonical section $\delta(\NL_{Y/X})$ to $\delta(\NL_{Y'/ X'})$, see
Derived Categories of Schemes, Remark \ref{perfect-remark-functorial-det}.

\begin{remark}
\label{remark-local-description-delta}
Let $Y \to X$ be a locally quasi-finite syntomic morphism of schemes.
What does the pair $(\det(\NL_{Y/X}), \delta(\NL_{Y/X}))$ look
like locally? Choose affine opens $V = \Spec(B) \subset Y$,
$U = \Spec(A) \subset X$ with $f(V) \subset U$ and an integer $n$ and
$f_1, \ldots, f_n \in A[x_1, \ldots, x_n]$ such that
$B = A[x_1, \ldots, x_n]/(f_1, \ldots, f_n)$. Then
$$
\NL_{B/A} = \left(
(f_1, \ldots, f_n)/(f_1, \ldots, f_n)^2
\longrightarrow
\bigoplus\nolimits_{i = 1, \ldots, n} B \text{d} x_i\right)
$$
and $(f_1, \ldots, f_n)/(f_1, \ldots, f_n)^2$ is free with generators
the classes $\overline{f}_i$. See proof of
Lemma \ref{lemma-syntomic-quasi-finite}.
Thus $\det(L_{B/A})$ is free on the generator
$$
\text{d}x_1 \wedge \ldots \wedge \text{d}x_n
\otimes
(\overline{f}_1 \wedge \ldots \wedge \overline{f}_n)^{\otimes -1}
$$
and the section $\delta(\NL_{B/A})$ is the element
$$
\delta(\NL_{B/A}) =
\det(\partial f_j/ \partial x_i) \cdot
\text{d}x_1 \wedge \ldots \wedge \text{d}x_n
\otimes
(\overline{f}_1 \wedge \ldots \wedge \overline{f}_n)^{\otimes -1}
$$
by definition.
\end{remark}

\noindent
Let $Y \to X$ be a locally quasi-finite syntomic morphism of
locally Noetherian schemes. By
Remarks \ref{remark-relative-dualizing-for-quasi-finite} and
\ref{remark-relative-dualizing-for-flat-quasi-finite} we have
a coherent $\mathcal{O}_Y$-module $\omega_{Y/X}$ and a canonical
global section
$$
\tau_{Y/X} \in \Gamma(Y, \omega_{Y/X})
$$
which affine locally recovers the pair $\omega_{B/A}, \tau_{B/A}$.
By Lemma \ref{lemma-dualizing-syntomic-quasi-finite} the module
$\omega_{Y/X}$ is invertible. Suppose given a commutative diagram of
locally Noetherian schemes
$$
\xymatrix{
Y' \ar[r]_b \ar[d] & Y \ar[d] \\
X' \ar[r] & X
}
$$
whose vertical arrows are locally quasi-finite syntomic and which
induces an isomorphism of $Y'$ with an open of $X' \times_X Y$.
Then there is a canonical base change map
$$
b^*\omega_{Y/X} \longrightarrow \omega_{Y'/X'}
$$
which is an isomorphism
mapping $\tau_{Y/X}$ to $\tau_{Y'/X'}$. Namely, the base change map
in the affine setting is (\ref{equation-bc-dualizing}), it is an
isomorphism by Lemma \ref{lemma-dualizing-base-change-of-flat}, and it
maps $\tau_{Y/X}$ to $\tau_{Y'/X'}$ by
Lemma \ref{lemma-trace-base-change} part (1).

\begin{lemma}
\label{lemma-syntomic-quasi-finite}
Let $f : Y \to X$ be a morphism of schemes. The following are equivalent
\begin{enumerate}
\item $f$ is locally quasi-finite and syntomic,
\item $f$ is locally quasi-finite, flat, and a local complete intersection
morphism,
\item $f$ is locally quasi-finite, flat, locally of finite presentation,
and the fibres of $f$ are local complete intersections,
\item $f$ is locally quasi-finite and for every $y \in Y$ there are
affine opens $y \in V = \Spec(B) \subset Y$, $U = \Spec(A) \subset X$
with $f(V) \subset U$ an integer $n$ and
$h, f_1, \ldots, f_n \in A[x_1, \ldots, x_n]$ such that
$B = A[x_1, \ldots, x_n, 1/h]/(f_1, \ldots, f_n)$,
\item for every $y \in Y$ there are affine opens
$y \in V = \Spec(B) \subset Y$, $U = \Spec(A) \subset X$
with $f(V) \subset U$ such that $A \to B$ is a relative global complete
intersection of the form $B = A[x_1, \ldots, x_n]/(f_1, \ldots, f_n)$,
\item $f$ is locally quasi-finite, flat, locally of finite presentation,
and $\NL_{Y/X}$ has tor-amplitude in $[-1, 0]$, and
\item $f$ is flat, locally of finite presentation,
$\NL_{Y/X}$ is perfect of rank $0$ with tor-amplitude in $[-1, 0]$,
\end{enumerate}
\end{lemma}

\begin{proof}
The equivalence of (1) and (2) is
More on Morphisms, Lemma \ref{more-morphisms-lemma-flat-lci}.
The equivalence of (1) and (3) is 
Morphisms, Lemma \ref{morphisms-lemma-syntomic-flat-fibres}.

\medskip\noindent
If $A \to B$ is as in (4), then
$B = A[x, x_1, \ldots, x_n]/(xh - 1, f_1, \ldots, f_n)$
is a relative global complete intersection by see Algebra, Definition
\ref{algebra-definition-relative-global-complete-intersection}.
Thus (4) implies (5).
It is clear that (5) implies (4).

\medskip\noindent
Condition (5) implies (1): by
Algebra, Lemma \ref{algebra-lemma-relative-global-complete-intersection}
a relative global complete intersection is syntomic and
the definition of a relative global complete intersection
guarantees that a relative global complete intersection on
$n$ variables with $n$ equations is quasi-finite, see
Algebra, Definition
\ref{algebra-definition-relative-global-complete-intersection} and
Lemma \ref{algebra-lemma-isolated-point-fibre}.

\medskip\noindent
Either Algebra, Lemma \ref{algebra-lemma-syntomic} or
Morphisms, Lemma \ref{morphisms-lemma-syntomic-locally-standard-syntomic}
shows that (1) implies (5).

\medskip\noindent
More on Morphisms, Lemma \ref{more-morphisms-lemma-flat-fp-NL-lci} shows that
(6) is equivalent to (1). If the equivalent conditions (1) -- (6) hold,
then we see that affine locally $Y \to X$ is given by a relative global
complete intersection $B = A[x_1, \ldots, x_n]/(f_1, \ldots, f_n)$
with the same number of variables as the number of
equations. Using this presentation we see that
$$
\NL_{B/A} =\left(
(f_1, \ldots, f_n)/(f_1, \ldots, f_n)^2
\longrightarrow
\bigoplus\nolimits_{i = 1, \ldots, n} B \text{d} x_i\right)
$$
By Algebra, Lemma
\ref{algebra-lemma-relative-global-complete-intersection-conormal}
the module $(f_1, \ldots, f_n)/(f_1, \ldots, f_n)^2$
is free with generators the congruence classes of the elements
$f_1, \ldots, f_n$. Thus $\NL_{B/A}$ has rank $0$ and so does $\NL_{Y/X}$.
In this way we see that (1) -- (6) imply (7).

\medskip\noindent
Finally, assume (7). By
More on Morphisms, Lemma \ref{more-morphisms-lemma-flat-fp-NL-lci}
we see that $f$ is syntomic. Thus on suitable affine opens
$f$ is given by a relative global complete intersection
$A \to B = A[x_1, \ldots, x_n]/(f_1, \ldots, f_m)$, see
Morphisms, Lemma \ref{morphisms-lemma-syntomic-locally-standard-syntomic}.
Exactly as above we see that $\NL_{B/A}$ is a perfect complex
of rank $n - m$. Thus $n = m$ and we see that (5) holds.
This finishes the proof.
\end{proof}

\begin{lemma}
\label{lemma-characterize-invertible}
Invertibility of the relative dualizing module.
\begin{enumerate}
\item If $A \to B$ is a quasi-finite flat homomorphism of Noetherian rings,
then $\omega_{B/A}$ is an invertible $B$-module if and only if
$\omega_{B \otimes_A \kappa(\mathfrak p)/\kappa(\mathfrak p)}$
is an invertible $B \otimes_A \kappa(\mathfrak p)$-module
for all primes $\mathfrak p \subset A$.
\item If $Y \to X$ is a quasi-finite flat morphism of
Noetherian schemes, then $\omega_{Y/X}$ is invertible
if and only if $\omega_{Y_x/x}$ is invertible for all $x \in X$.
\end{enumerate}
\end{lemma}

\begin{proof}
Proof of (1). As $A \to B$ is flat, the module
$\omega_{B/A}$ is $A$-flat, see Lemma \ref{lemma-dualizing-base-flat-flat}.
Thus $\omega_{B/A}$ is an invertible $B$-module if and only if
$\omega_{B/A} \otimes_A \kappa(\mathfrak p)$
is an invertible $B \otimes_A \kappa(\mathfrak p)$-module for
every prime $\mathfrak p \subset A$, see More on Morphisms, Lemma
\ref{more-morphisms-lemma-flat-and-free-at-point-fibre}.
Still using that $A \to B$ is flat, we have that
formation of $\omega_{B/A}$ commutes with base change, see
Lemma \ref{lemma-dualizing-base-change-of-flat}.
Thus we see that invertibility of the relative dualizing module,
in the presence of flatness, is equivalent to invertibility
of the relative dualizing module for the maps
$\kappa(\mathfrak p) \to B \otimes_A \kappa(\mathfrak p)$.

\medskip\noindent
Part (2) follows from (1) and the fact that affine locally
the dualizing modules are given by their algebraic counterparts, see
Remark \ref{remark-relative-dualizing-for-quasi-finite}.
\end{proof}

\begin{lemma}
\label{lemma-dim-zero-global-complete-intersection-over-field}
Let $k$ be a field. Let $B = k[x_1, \ldots, x_n]/(f_1, \ldots, f_n)$
be a global complete intersection over $k$ of dimension $0$.
Then $\omega_{B/k}$ is invertible.
\end{lemma}

\begin{proof}
By Noether normalization, see
Algebra, Lemma \ref{algebra-lemma-Noether-normalization}
we see that there exists a finite injection $k \to B$, i.e.,
$\dim_k(B) < \infty$. Hence $\omega_{B/k} = \Hom_k(B, k)$
as a $B$-module.
By Dualizing Complexes, Lemma \ref{dualizing-lemma-dualizing-finite}
we see that $R\Hom(B, k)$ is a dualizing complex for $B$
and by Dualizing Complexes, Lemma \ref{dualizing-lemma-RHom-ext}
we see that $R\Hom(B, k)$ is equal to $\omega_{B/k}$
placed in degree $0$. Thus it suffices to show that
$B$ is Gorenstein
(Dualizing Complexes, Lemma \ref{dualizing-lemma-gorenstein}).
This is true by Dualizing Complexes, Lemma
\ref{dualizing-lemma-gorenstein-lci}.
\end{proof}

\begin{lemma}
\label{lemma-dualizing-syntomic-quasi-finite}
Let $f : Y \to X$ be a morphism of locally Noetherian schemes. If $f$
satisfies the equivalent conditions of Lemma \ref{lemma-syntomic-quasi-finite}
then $\omega_{Y/X}$ is an invertible $\mathcal{O}_Y$-module.
\end{lemma}

\begin{proof}
We may assume $A \to B$ is a relative global complete
intersection of the form $B = A[x_1, \ldots, x_n]/(f_1, \ldots, f_n)$
and we have to show $\omega_{B/A}$ is invertible.
This follows in combining Lemmas \ref{lemma-characterize-invertible} and
\ref{lemma-dim-zero-global-complete-intersection-over-field}.
\end{proof}

\begin{example}
\label{example-universal-quasi-finite-syntomic}
Let $n \geq 1$ and $d \geq 1$ be integers. Let $T$ be the set of
multi-indices $E = (e_1, \ldots, e_n)$ with $e_i \geq 0$ and
$\sum e_i \leq d$. Consider the ring
$$
A = \mathbf{Z}[a_{i, E} ; 1 \leq i \leq n, E \in T]
$$
In $A[x_1, \ldots, x_n]$ consider the elements
$f_i = \sum_{E \in T} a_{i, E} x^E$ where $x^E = x_1^{e_1} \ldots x_n^{e_n}$
as is customary. Consider the $A$-algebra
$$
B = A[x_1, \ldots, x_n]/(f_1, \ldots, f_n)
$$
Denote $X_{n, d} = \Spec(A)$ and let $Y_{n, d} \subset \Spec(B)$
be the maximal open subscheme such that the restriction of the
morphism $\Spec(B) \to \Spec(A) = X_{n, d}$ is quasi-finite, see
Algebra, Lemma \ref{algebra-lemma-quasi-finite-open}.
\end{example}

\begin{lemma}
\label{lemma-universal-quasi-finite-syntomic-etale}
With notation as in Example \ref{example-universal-quasi-finite-syntomic}
the schemes $X_{n, d}$ and $Y_{n, d}$ are regular and irreducible,
the morphism $Y_{n, d} \to X_{n, d}$ is locally quasi-finite and
syntomic, and there is a dense open subscheme $V \subset Y_{n, d}$
such that $Y_{n, d} \to X_{n, d}$ restricts to an \'etale morphism
$V \to X_{n, d}$.
\end{lemma}

\begin{proof}
The scheme $X_{n, d}$ is the spectrum of the polynomial ring $A$.
Hence $X_{n, d}$ is regular and irreducible. Since we can write
$$
f_i = a_{i, (0, \ldots, 0)} +
\sum\nolimits_{E \in T, E \not = (0, \ldots, 0)} a_{i, E} x^E
$$
we see that the ring $B$ is isomorphic to the polynomial ring
on $x_1, \ldots, x_n$ and the elements $a_{i, E}$ with
$E \not = (0, \ldots, 0)$. Hence $\Spec(B)$ is an irreducible and
regular scheme and so is the open $Y_{n, d}$. The morphism
$Y_{n, d} \to X_{n, d}$ is locally quasi-finite and syntomic by
Lemma \ref{lemma-syntomic-quasi-finite}. To find $V$ it suffices
to find a single point where $Y_{n, d} \to X_{n, d}$ is \'etale
(the locus of points where a morphism is \'etale is open by
definition). Thus it suffices to find a point of $X_{n, d}$
where the fibre of $Y_{n, d} \to X_{n, d}$ is nonempty and \'etale, see
Morphisms, Lemma \ref{morphisms-lemma-etale-at-point}. We choose
the point corresponding to the ring map $\chi : A \to \mathbf{Q}$
sending $f_i$ to $x_i^d - 1$. Then
$$
B \otimes_{A, \chi} \mathbf{Q} =
\mathbf{Q}[x_1, \ldots, x_n]/(x_1^d - 1, \ldots, x_n^d - 1)
$$
which is a nonzero \'etale algebra over $\mathbf{Q}$.
\end{proof}

\begin{lemma}
\label{lemma-locally-comes-from-universal}
Let $f : Y \to X$ be a morphism of schemes. If $f$ satisfies the equivalent
conditions of Lemma \ref{lemma-syntomic-quasi-finite} then for every
$y \in Y$ there exist $n, d$ and a commutative diagram
$$
\xymatrix{
Y \ar[d] &
V \ar[d] \ar[l] \ar[r] &
Y_{n, d} \ar[d] \\
X & U \ar[l] \ar[r] &
X_{n, d}
}
$$
where $U \subset X$ and $V \subset Y$ are open with $y \in V$,
where $Y_{n, d} \to X_{n, d}$
is as in Example \ref{example-universal-quasi-finite-syntomic}, and
where the square on the right hand side is cartesian.
\end{lemma}

\begin{proof}
By Lemma \ref{lemma-syntomic-quasi-finite}
we can choose $U$ and $V$ affine so that
$U = \Spec(R)$ and $V = \Spec(S)$ with
$S = R[y_1, \ldots, y_n]/(g_1, \ldots, g_n)$.
With notation as in Example \ref{example-universal-quasi-finite-syntomic}
if we pick $d$ large enough, then we can write each $g_i$ as
$g_i = \sum_{E \in T} g_{i, E}y^E$ with $g_{i, E} \in R$.
Then the map $A \to R$ sending $a_{i, E}$ to $g_{i, E}$
and the map $B \to S$ sending $x_i \to y_i$ give a cocartesian
diagram of rings
$$
\xymatrix{
S & B \ar[l] \\
R \ar[u] & A \ar[l] \ar[u]
}
$$
which proves the lemma.
\end{proof}

\section{A formula for the different}
\label{section-formula-different}

\noindent
In this section we discuss the material in \cite[Appendix A]{Mazur-Roberts}
due to Tate. In our language, this will show that the different is
equal to the K\"ahler different in the case of a flat, quasi-finite,
local complete intersection morphism.
First we compute the Noether different in a special case.

\begin{lemma}
\label{lemma-tate}
\begin{reference}
\cite[Appendix]{Mazur-Roberts}
\end{reference}
Let $A \to P$ be a ring map. Let $f_1, \ldots, f_n \in P$ be a
Koszul regular sequence. Assume $B = P/(f_1, \ldots, f_n)$
is flat over $A$. Let $g_1, \ldots, g_n \in P \otimes_A B$
be a Koszul regular sequence generating the kernel of the multiplication
map $P \otimes_A B \to B$. Write $f_i \otimes 1 = \sum g_{ij} g_j$.
Then the annihilator of $\Ker(B \otimes_A B \to B)$ is a principal
ideal generated by the image of $\det(g_{ij})$.
\end{lemma}

\begin{proof}
The Koszul complex $K_\bullet = K(P, f_1, \ldots, f_n)$ is a resolution
of $B$ by finite free $P$-modules. The Koszul complex
$M_\bullet = K(P \otimes_A B, g_1, \ldots, g_n)$ is a resolution
of $B$ by finite free $P \otimes_A B$-modules. There is a map of
complexes
$$
K_\bullet \longrightarrow M_\bullet
$$
which in degree $1$ is given by the matrix $(g_{ij})$ and
in degree $n$ by $\det(g_{ij})$. See
More on Algebra, Lemma \ref{more-algebra-lemma-functorial}.
As $B$ is a flat $A$-module, we can view $M_\bullet$ as a complex
of flat $P$-modules (via $P \to P \otimes_A B$, $p \mapsto p \otimes 1$).
Thus we may use both complexes to compute $\text{Tor}_*^P(B, B)$ and
it follows that the displayed map defines a quasi-isomorphism after tensoring
with $B$. It is clear that $H_n(K_\bullet \otimes_P B) = B$.
On the other hand, $H_n(M_\bullet \otimes_P B)$ is the kernel of
$$
B \otimes_A B \xrightarrow{g_1, \ldots, g_n} (B \otimes_A B)^{\oplus n}
$$
Since $g_1, \ldots, g_n$ generate the kernel of $B \otimes_A B \to B$
this proves the lemma.
\end{proof}

\begin{lemma}
\label{lemma-quasi-finite-complete-intersection}
Let $A$ be a ring. Let $n \geq 1$ and
$h, f_1, \ldots, f_n \in A[x_1, \ldots, x_n]$.
Set $B = A[x_1, \ldots, x_n, 1/h]/(f_1, \ldots, f_n)$.
Assume that $B$ is quasi-finite over $A$.
Then
\begin{enumerate}
\item $B$ is flat over $A$ and $A \to B$ is a relative local complete
intersection,
\item the annihilator $J$ of $I = \Ker(B \otimes_A B \to B)$
is free of rank $1$ over $B$,
\item the Noether different of $B$ over $A$ is generated
by $\det(\partial f_i/\partial x_j)$ in $B$.
\end{enumerate}
\end{lemma}

\begin{proof}
Note that
$B = A[x, x_1, \ldots, x_n]/(xh - 1, f_1, \ldots, f_n)$
is a relative global complete intersection over $A$, see
Algebra, Definition
\ref{algebra-definition-relative-global-complete-intersection}.
By Algebra, Lemma \ref{algebra-lemma-relative-global-complete-intersection}
we see that $B$ is flat over $A$.

\medskip\noindent
Write $P' = A[x, x_1, \ldots, x_n]$ and
$P = P'/(xh - 1) = A[x_1, \ldots, x_n, 1/h]$.
Then we have $P' \to P \to B$.
By More on Algebra, Lemma
\ref{more-algebra-lemma-relative-global-complete-intersection-koszul}
we see that $xh - 1, f_1, \ldots, f_n$ is a Koszul regular sequence
in $P'$. Since $xh - 1$ is a Koszul regular sequence of length
one in $P'$ (by the same lemma for example) we conclude that 
$f_1, \ldots, f_n$ is a Koszul regular sequence in $P$ by
More on Algebra, Lemma \ref{more-algebra-lemma-truncate-koszul-regular}.

\medskip\noindent
Let $g_i \in P \otimes_A B$ be the image of $x_i \otimes 1 - 1 \otimes x_i$.
Let us use the short hand $y_i = x_i \otimes 1$ and $z_i = 1 \otimes x_i$
in $A[x_1, \ldots, x_n] \otimes_A A[x_1, \ldots, x_n]$
so that $g_i$ is the image of $y_i - z_i$. For a polynomial
$f \in A[x_1, \ldots, x_n]$ we write $f(y) = f \otimes 1$
and $f(z) = 1 \otimes f$ in the above tensor product.
Then we have
$$
P \otimes_A B/(g_1, \ldots, g_n) =
\frac{A[y_1, \ldots, y_n, z_1, \ldots, z_n, \frac{1}{h(y)h(z)}]}
{(f_1(z), \ldots, f_n(z), y_1 - z_1, \ldots, y_n - z_n)}
$$
which is clearly isomorphic to $B$. Hence by the same arguments
as above we find that $f_1(z), \ldots, f_n(z), y_1 - z_1, \ldots, y_n - z_n$
is a Koszul regular sequence in
$A[y_1, \ldots, y_n, z_1, \ldots, z_n, \frac{1}{h(y)h(z)}]$.
The sequence $f_1(z), \ldots, f_n(z)$ is a Koszul regular in
$A[y_1, \ldots, y_n, z_1, \ldots, z_n, \frac{1}{h(y)h(z)}]$
by flatness of the map
$$
P \longrightarrow A[y_1, \ldots, y_n, z_1, \ldots, z_n,
\textstyle{\frac{1}{h(y)h(z)}}],\quad x_i \longmapsto z_i
$$
and More on Algebra, Lemma
\ref{more-algebra-lemma-koszul-regular-flat-base-change}.
By More on Algebra, Lemma \ref{more-algebra-lemma-truncate-koszul-regular}
we conclude that $g_1, \ldots, g_n$ is a regular sequence
in $P \otimes_A B$.

\medskip\noindent
At this point we have verified all the assumptions of Lemma \ref{lemma-tate}
above with $P$, $f_1, \ldots, f_n$, and $g_i \in P \otimes_A B$ as above.
In particular the annihilator $J$ of $I$ is freely generated by one
element $\delta$ over $B$.
Set $f_{ij} = \partial f_i/\partial x_j \in A[x_1, \ldots, x_n]$.
An elementary computation shows that we can write
$$
f_i(y) =
f_i(z_1 + g_1, \ldots, z_n + g_n) =
f_i(z) + \sum\nolimits_j f_{ij}(z) g_j +
\sum\nolimits_{j, j'} F_{ijj'}g_jg_{j'}
$$
for some $F_{ijj'} \in A[y_1, \ldots, y_n, z_1, \ldots, z_n]$.
Taking the image in $P \otimes_A B$ the terms $f_i(z)$ map to
zero and we obtain
$$
f_i \otimes 1 = \sum\nolimits_j
\left(1 \otimes f_{ij} + \sum\nolimits_{j'} F_{ijj'}g_{j'}\right)g_j
$$
Thus we conclude from Lemma \ref{lemma-tate}
that $\delta = \det(g_{ij})$ with
$g_{ij} = 1 \otimes f_{ij} + \sum_{j'} F_{ijj'}g_{j'}$.
Since $g_{j'}$ maps to zero in $B$, we conclude
that the image of $\det(\partial f_i/\partial x_j)$ in $B$
generates the Noether different of $B$ over $A$.
\end{proof}

\begin{lemma}
\label{lemma-different-syntomic-quasi-finite}
Let $f : Y \to X$ be a morphism of Noetherian schemes. If $f$
satisfies the equivalent conditions of Lemma \ref{lemma-syntomic-quasi-finite}
then the different $\mathfrak{D}_f$ of $f$ is the K\"ahler different
of $f$.
\end{lemma}

\begin{proof}
By Lemmas \ref{lemma-flat-gorenstein-agree-noether} and
\ref{lemma-dualizing-syntomic-quasi-finite}
the different of $f$ affine locally is the same as the
Noether different. Then the lemma follows from the
computation of the Noether different and the K\"ahler
different on standard affine pieces done in
Lemmas \ref{lemma-kahler-different-complete-intersection} and
\ref{lemma-quasi-finite-complete-intersection}.
\end{proof}

\begin{lemma}
\label{lemma-different-quasi-finite-complete-intersection}
Let $A$ be a ring. Let $n \geq 1$ and
$h, f_1, \ldots, f_n \in A[x_1, \ldots, x_n]$.
Set $B = A[x_1, \ldots, x_n, 1/h]/(f_1, \ldots, f_n)$.
Assume that $B$ is quasi-finite over $A$.
Then there is an isomorphism $B \to \omega_{B/A}$
mapping $\det(\partial f_i/\partial x_j)$ to $\tau_{B/A}$.
\end{lemma}

\begin{proof}
Let $J$ be the annihilator of $\Ker(B \otimes_A B \to B)$.
By Lemma \ref{lemma-quasi-finite-complete-intersection}
the map $A \to B$ is flat and
$J$ is a free $B$-module with generator $\xi$ mapping to
$\det(\partial f_i/\partial x_j)$ in $B$.
Thus the lemma follows from
Lemma \ref{lemma-noether-different-flat-quasi-finite}
and the fact (Lemma \ref{lemma-dualizing-syntomic-quasi-finite})
that $\omega_{B/A}$ is an invertible $B$-module.
(Warning: it is necessary to prove $\omega_{B/A}$
is invertible because a finite $B$-module $M$ such
that $\Hom_B(M, B) \cong B$ need not be free.)
\end{proof}

\begin{proof}
Let us reformulate the statement of the proposition. Consider the category
$\mathcal{C}$ whose objects, denoted $Y/X$, are locally quasi-finite syntomic
morphism $Y \to X$ of locally Noetherian schemes and whose morphisms
$b/a : Y'/X' \to Y/X$ are commutative diagrams
$$
\xymatrix{
Y' \ar[d] \ar[r]_b & Y \ar[d] \\
X' \ar[r]^a & X
}
$$
which induce an isomorphism of $Y'$ with an open subscheme of
$X' \times_X Y$. The proposition means that for every object
$Y/X$ of $\mathcal{C}$ we have an isomorphism
$c_{Y/X} : \det(\NL_{Y/X}) \to \omega_{Y/X}$
with $c_{Y/X}(\delta(\NL_{Y/X})) = \tau_{Y/X}$
and for every morphism $b/a : Y'/X' \to Y/X$ of $\mathcal{C}$ we have
$b^*c_{Y/X} = c_{Y'/X'}$ via the identifications
$b^*\det(\NL_{Y/X}) = \det(\NL_{Y'/X'})$ and
$b^*\omega_{Y/X} = \omega_{Y'/X'}$ described above.

\medskip\noindent
Given $Y/X$ in $\mathcal{C}$ and $y \in Y$ we can find
an affine opens $V \subset Y$ and $U \subset X$ with
$y \in V$ and $f(V) \subset U$
such that there exists some isomorphism
$$
\det(\NL_{Y/X})|_V \longrightarrow \omega_{Y/X}|_V
$$
mapping $\delta(\NL_{Y/X})|_V$ to $\tau_{Y/X}|_V$. This follows
from picking affine opens as in
Lemma \ref{lemma-syntomic-quasi-finite} part (5), the affine
local description of $\delta(\NL_{Y/X})$ in
Remark \ref{remark-local-description-delta}, and
Lemma \ref{lemma-different-quasi-finite-complete-intersection}.
If the annihilator of the section $\tau_{Y/X}$ is zero, then
these local maps are unique and automatically glue. Hence if the annihilator
of $\tau_{Y/X}$ is zero, then there is a unique isomorphism
$c_{Y/X} : \det(\NL_{Y/X}) \to \omega_{Y/X}$ with
$c_{Y/X}(\delta(\NL_{Y/X})) = \tau_{Y/X}$.
If $b/a : Y'/X' \to Y/X$ is a morphism of $\mathcal{C}$
and the annihilator of $\tau_{Y'/X'}$ is zero as well,
then $b^*c_{Y/X}$ is the unique isomorphism
$c_{Y'/X'} : \det(\NL_{Y'/X'}) \to \omega_{Y'/X'}$ with
$c_{Y'/X'}(\delta(\NL_{Y'/X'})) = \tau_{Y'/X'}$.
This follows formally from the fact that
$b^*\delta(\NL_{Y/X}) = \delta(\NL_{Y'/X'})$ and
$b^*\tau_{Y/X} = \tau_{Y'/X'}$.

\medskip\noindent
We can summarize the results of the previous paragraph as follows.
Let $\mathcal{C}_{nice} \subset \mathcal{C}$ denote the
full subcategory of $Y/X$ such that the annihilator of
$\tau_{Y/X}$ is zero. Then we have solved the problem
on $\mathcal{C}_{nice}$. For $Y/X$ in $\mathcal{C}_{nice}$
we continue to denote $c_{Y/X}$ the solution we've just found.

\medskip\noindent
Consider morphisms
$$
Y_1/X_1 \xleftarrow{b_1/a_1} Y/X \xrightarrow{b_2/a_2} Y_2/X_2
$$
in $\mathcal{C}$ such that $Y_1/X_1$ and $Y_2/X_2$ are objects
of $\mathcal{C}_{nice}$. {\bf Claim.} $b_1^*c_{Y_1/X_1} = b_2^*c_{Y_2/X_2}$.
We will first show that the claim implies the proposition
and then we will prove the claim.

\medskip\noindent
Let $d, n \geq 1$ and consider the locally
quasi-finite syntomic morphism $Y_{n, d} \to X_{n, d}$
constructed in Example \ref{example-universal-quasi-finite-syntomic}.
Then $Y_{n, d}$ is an irreducible regular scheme and the
morphism $Y_{n, d} \to X_{n, d}$ is locally quasi-finite syntomic
and \'etale over a dense open, see
Lemma \ref{lemma-universal-quasi-finite-syntomic-etale}.
Thus $\tau_{Y_{n, d}/X_{n, d}}$ is nonzero for example by
Lemma \ref{lemma-different-ramification}. Now a nonzero section
of an invertible module over an irreducible regular scheme
has vanishing annihilator. Thus
$Y_{n, d}/X_{n, d}$ is an object of $\mathcal{C}_{nice}$.

\medskip\noindent
Let $Y/X$ be an arbitrary object of $\mathcal{C}$. Let $y \in Y$.
By Lemma \ref{lemma-locally-comes-from-universal} we can find
$n, d \geq 1$ and morphisms
$$
Y/X \leftarrow V/U \xrightarrow{b/a} Y_{n, d}/X_{n, d}
$$
of $\mathcal{C}$ such that $V \subset Y$ and $U \subset X$ are open.
Thus we can pullback the canonical morphism $c_{Y_{n, d}/X_{n, d}}$
constructed above by $b$ to $V$. The claim guarantees these local
isomorphisms glue! Thus we get a well defined global isomorphism
$c_{Y/X} : \det(\NL_{Y/X}) \to \omega_{Y/X}$ with
$c_{Y/X}(\delta(\NL_{Y/X})) = \tau_{Y/X}$.
If $b/a : Y'/X' \to Y/X$ is a morphism of $\mathcal{C}$, then
the claim also implies that the similarly constructed map
$c_{Y'/X'}$ is the pullback by $b$ of the locally constructed
map $c_{Y/X}$. Thus it remains to prove the claim.

\medskip\noindent
In the rest of the proof we prove the claim. We may pick a point
$y \in Y$ and prove the maps agree in an open neighbourhood of $y$.
Thus we may replace $Y_1$, $Y_2$ by open neighbourhoods of the
image of $y$ in $Y_1$ and $Y_2$. Therefore we may assume
$Y, X, Y_1, X_1, Y_2, X_2$ are affine, say they are the spectra
of rings $B, A, B_1, A_1, B_2, A_2$. Picture
$$
\xymatrix{
B_1 \ar[r] & B & B_2 \ar[l] \\
A_1 \ar[u] \ar[r] & A \ar[u] & A_2 \ar[l] \ar[u]
}
$$
By assumption the spectrum of $B$ is an affine open of
both the spectrum of $A \otimes_{A_1} B_1$ and $A \otimes_{A_2} B_2$.
Shrinking more we may assume there exist elements
$g_i \in A \otimes_{A_i} B_i$ such that our maps give isomorphisms
$(A \otimes_{A_i} B_i)_{g_i} = B$, see
Properties, Lemma \ref{properties-lemma-standard-open-two-affines}.
Let $x_\alpha$, $y_\beta$ be a sufficiently large collection of
variables such that we may choose surjections
$$
A'_1 = A_1[x_\alpha] \to A
\quad\text{and}\quad
A'_2 = A_2[y_\beta] \to A
$$
of $A_1$ and $A_2$-algebras. Then we can choose lifts
$h_i \in A'_i \otimes_{A_i} B_i$ of $g_i \in A \otimes_{A_i} B_i$
and we consider the diagram
$$
\xymatrix{
(A'_1 \otimes_{A_1} B_1)_{h_1} \ar[r] & B &
(A'_2 \otimes_{A_1} B_2)_{h_2} \ar[l] \\
A'_1 \ar[u] \ar[r] & A \ar[u] & A'_2 \ar[l] \ar[u]
}
$$
By construction the two squares are cocartesian. Next, we consider the
ring map
$$
A' = A'_1 \times_A A'_2 \longrightarrow
B' =
(A'_1 \otimes_{A_1} B_1)_{h_1} \times_B
(A'_2 \otimes_{A_1} B_2)_{h_2}
$$
By
More on Algebra, Lemma \ref{more-algebra-lemma-module-over-fibre-product}
we have $A'_1 \otimes_{A'} B' = (A'_1 \otimes_{A_1} B_1)_{h_1}$
and $A'_2 \otimes_{A'} B' = (A'_2 \otimes_{A_2} B_2)_{h_2}$. In particular
the fibres of the morphism $Y' = \Spec(B') \to \Spec(A') = X'$ are
open subschemes of base changes of the fibres of the maps $Y_i \to X_i$,
hence finite and local complete intersections
(see Lemma \ref{lemma-syntomic-quasi-finite}).
By More on Algebra, Lemma
\ref{more-algebra-lemma-properties-algebras-over-fibre-product}
the ring map $A' \to B'$ is flat and of finite presentation.
Thus by Lemma \ref{lemma-syntomic-quasi-finite} part (3)
we conclude that $Y' \to X'$ is locally quasi-finite and syntomic.
Thus we obtain the commutative diagram
$$
\xymatrix{
& & Y/X \ar[ld] \ar@/_2pc/[lldd]_{b_1/a_1} \ar[rd] \ar@/^2pc/[rrdd]^{b_2/a_2} \\
& Y'_1/X'_1 \ar[rd] \ar[ld] & & Y'_2/X'_2 \ar[ld] \ar[rd] \\
Y_1/X_1 & & Y'/X' & & Y_2/X_2
}
$$
where $Y'_i/X'_i$ corresponds to $A'_i \to (A'_i \otimes_{A_i} B_i)_{h_i}$.

\medskip\noindent
In this paragraph, we fix the issue that $A'$ and $A'_i$ may not be Noetherian
by a limit argument; we suggest the reader skip this.
We can find a finitely generated $\mathbf{Z}$-subalgebra
$A'' \subset A'$ and a quasi-finite and syntomic ring map
$A'' \to B''$ such that $B' = A' \otimes_{A''} B''$.
This follows from Limits, Lemmas
\ref{limits-lemma-descend-finite-presentation},
\ref{limits-lemma-descend-syntomic}, and
\ref{limits-lemma-descend-quasi-finite}.
Then the ring maps $A'' \to A'_1 = A_1[x_\alpha]$ will factor through a finite
polynomial subalgebra $A_i[x_1, \ldots, x_n$ and the ring map
$A'' \to A'_2 = A_2[y_\beta]$ will factor through $A_2[y_1, \ldots, y_m]$.
After increasing the sets of variables, we may also assume that
$h_1 \in A_1[x_1, \ldots, x_n] \otimes_{A_1} B_1$ and
$h_2 \in A_2[y_1, \ldots, y_m] \otimes_{A_2} B_2$.
Increasing the sets of variables one more time, we can assume
that the maps $B'' \to (A'_i \otimes_{A_i} B_i)_{h_i}$
factor through $(A_1[x_1, \ldots, x_n] \otimes_{A_1} B_1)_{h_1}$
and $(A_2[y_1, \ldots, y_m] \otimes_{A_2} B_2)_{h_2}$.
Replacing $Y'$, $X'$, $Y'_1$, $X'_1$, $Y'_2$, $X'_2$ by the spectra of
$B''$, $A''$, $(A_1[x_1, \ldots, x_n] \otimes_{A_1} B_1)_{h_1}$,
$A_1[x_1, \ldots, x_n]$, $(A_2[y_1, \ldots, y_m] \otimes_{A_2} B_2)_{h_2}$,
$A_2[y_1, \ldots, y_m]$, we see that we get a diagram as above
where all schemes are Noetherian.

\medskip\noindent
Note that $Y'_i/X'_i$ is an object of $\mathcal{C}_{nice}$ as it
obtained by taking an open subscheme of the product of an affine space
with $Y_i/X_i$. In particular, the pullback of
$c_{Y_i/X_i}$ via $(b_i/a_i)$ is the same as the pullback of
$c_{Y'_i/X'_i}$ via $Y/X \to Y'_i/X'_i$.
Now we would be done if $Y'/X'$ is an object of $\mathcal{C}_{nice}$
(but this is likely not the case). Namely, then pulling back $c_{Y'/X'}$
around the two sides of the square, we would obtain the desired conclusion.
To get around the problem that $Y'/X'$ is not in $\mathcal{C}_{nice}$
we note the arguments above show that, after possibly shrinking all
of the schemes $X, Y, X'_1, Y'_1, X'_2, Y'_2, X', Y'$ we can find some
$n, d \geq 1$, and extend the diagram like so:
$$
\xymatrix{
& Y/X \ar[ld] \ar[rd] \\
Y'_1/X'_1 \ar[rd] & & Y'_2/X'_2 \ar[ld] \\
& Y'/X' \ar[d] \\
& Y_{n, d}/X_{n, d}
}
$$
and then we can use the already given argument by pulling
back from $c_{Y_{n, d}/X_{n, d}}$. This finishes the proof.
\end{proof}

\bibliographystyle{plain}
\bibliography{my}
\end{document}
